\section{Problem Formulation and Recourse Structure}
\label{sec:model}

This section formulates the wireless charging infrastructure planning problem as a two-stage stochastic program.
The formulation separates long-lived wireless-facility decisions from scenario-adaptive battery and charging decisions, and then identifies the recourse structure used by the algorithm in Section~\ref{sec:method}.

\subsection{Two-stage stochastic formulation}
\label{subsec:two-stage-formulation}

Let $\Omega$ be the set of service-plan scenarios, and let $p_\omega$ be the probability of scenario $\omega \in \Omega$, where $p_\omega>0$ and $\sum_{\omega\in\Omega}p_\omega=1$.
Each scenario represents one plausible future service plan and fixes the route frequencies, timetables, bus trip chains, active-fleet requirements, and trip-level energy requirements, while the operated route set $\mathcal{R}$ is common across scenarios.
The planner first installs long-lived wireless charging facilities before the future service plan is known.
After a scenario is realized, the second-stage recourse selects route-level battery capacities and the corresponding representative-day charging allocation and energy-management decisions.

Let $\mathcal{S}$ be the station set, $\mathcal{L}$ the directed link set between consecutive stations, and $\mathcal{C}\subseteq\mathcal{S}$ the set of existing interchanges or hubs equipped with plug-in chargers.
Stationary wireless charging (SWC) can be installed at any station in $\mathcal{S}\setminus\mathcal{C}$, while dynamic wireless charging (DWC) can be installed on any directed link in $\mathcal{L}$.
The first-stage binary variables $w_s^{\mathrm{S}}$ and $w_l^{\mathrm{D}}$ indicate whether SWC is installed at station $s\in\mathcal{S}\setminus\mathcal{C}$ and DWC is installed on link $l\in\mathcal{L}$, respectively.
Existing plug-in chargers at stations in $\mathcal{C}$ are pre-existing resources with limited shared capacity and are not candidates for SWC deployment.

The compact two-stage stochastic formulation is
\begin{align}
    \min_{w^{\mathrm{S}},w^{\mathrm{D}}} \quad
        & C^{\mathrm{fac}}(w)
        +
        \sum_{\omega \in \Omega} p_\omega \Phi_\omega(w)
        \label{eq:two-stage-model} \\
    \mathrm{s.t.} \quad
        & w_s^{\mathrm{S}} \in \{0,1\},
        && s \in \mathcal{S}\setminus\mathcal{C},
        \label{eq:swc-domain} \\
        & w_l^{\mathrm{D}} \in \{0,1\},
        && l \in \mathcal{L}.
        \label{eq:dwc-domain}
\end{align}
where $w=(w^{\mathrm{S}},w^{\mathrm{D}})$ and
\[
    C^{\mathrm{fac}}(w)=\sum_{s \in \mathcal{S}\setminus\mathcal{C}} c_s^{\mathrm{S}} w_s^{\mathrm{S}}+\sum_{l \in \mathcal{L}} c_l^{\mathrm{D}} w_l^{\mathrm{D}}.
\]
Here $c_s^{\mathrm{S}}$ and $c_l^{\mathrm{D}}$ are the representative-day equivalent costs of installing SWC at station $s$ and DWC on link $l$, respectively.
Facility and battery investment costs are converted to representative-day equivalents to match the charging-cost scale.
The vector $w$ is the nonanticipative here-and-now facility decision, and $\Phi_\omega(w)$ is the optimal second-stage value in scenario $\omega$.

The scenario decisions consist of route-level battery selection variables $z_{\omega,r,k}$, plug-in charger occupancy variables collected in $x_\omega$, and remaining operating variables collected in $y_\omega$, including charging amounts and battery energy states.
For a fixed facility layout $w$, the scenario recourse value $\Phi_\omega(w)$ is the minimum scenario cost
\begin{equation}
    \Phi_\omega(w)
    =
    \min_{z_\omega,x_\omega,y_\omega}
    \left\{
        C^{\mathrm{bat}}_\omega(z_\omega) + C^{\mathrm{chg}}_\omega(y_\omega):
        (z_\omega,x_\omega,y_\omega)\in\mathcal{X}_\omega(w)
    \right\}.
    \label{eq:recourse-value}
\end{equation}
The next subsection expands $\mathcal{X}_\omega(w)$.

\subsection{Scenario recourse formulation}
\label{subsec:scenario-recourse-formulation}

For a fixed scenario $\omega$, the recourse formulation uses route-indexed trip-chain data.
For each route $r\in\mathcal R$, $\mathcal{M}_{\omega,r}$ is the set of trip chains, $\mathcal{A}_{\omega,r}$ is the set of directed traversals, and $\mathcal{V}^{\mathrm{P}}_{\omega,r}$ and $\mathcal{V}^{\mathrm{S}}_{\omega,r}$ are the plug-in and SWC station visits.
Each trip chain corresponds to one bus, so $n_{\omega,r}=|\mathcal{M}_{\omega,r}|$ is the route fleet size in scenario $\omega$.
A consolidated notation table is provided in Appendix~\ref{app:model-details}.

\paragraph{Recourse costs.}
Battery options $\mathcal{K}$ have capacities $\beta_k$, with representative-day equivalent unit cost $c^{\mathrm{bat}}$.
Operating cost is based on grid-side energy purchased through plug-in charging, stationary wireless charging, dynamic wireless charging, and overnight depot charging.
The corresponding charging amounts are $q_v^{\mathrm{P}}$, $q_v^{\mathrm{S}}$, $q_a^{\mathrm{D}}$, and $q_m^{\mathrm{N}}$, with nonnegative daytime and overnight electricity prices $c^{\mathrm{day}}$ and $c^{\mathrm{night}}$.
The scenario battery cost is
\[
    C^{\mathrm{bat}}_\omega(z_\omega)
    =
    c^{\mathrm{bat}}
    \sum_{r \in \mathcal{R}}
    n_{\omega,r}
    \sum_{k \in \mathcal{K}}
    \beta_k z_{\omega,r,k}.
\]
The scenario charging cost is
\[
    C^{\mathrm{chg}}_\omega(y_\omega)
    =
    c^{\mathrm{day}}
    \left(
    \sum_{r \in \mathcal{R}}
    \sum_{v \in \mathcal{V}^{\mathrm{P}}_{\omega,r}} q_v^{\mathrm{P}}
    +
    \sum_{r \in \mathcal{R}}
    \sum_{v \in \mathcal{V}^{\mathrm{S}}_{\omega,r}} q_v^{\mathrm{S}}
    +
    \sum_{r \in \mathcal{R}}
    \sum_{a \in \mathcal{A}_{\omega,r}} q_a^{\mathrm{D}}
    \right)
    +
    c^{\mathrm{night}}
    \sum_{r \in \mathcal{R}}
    \sum_{m \in \mathcal{M}_{\omega,r}}
    q_m^{\mathrm{N}}.
\]

\paragraph{Battery configuration.}
The route-level battery-choice constraints are
\begin{align}
    & \sum_{k \in \mathcal{K}} z_{\omega,r,k} = 1,
    && r \in \mathcal{R},
    \label{eq:one-battery} \\
    & B_{\omega,r}
    =
    \sum_{k \in \mathcal{K}} \beta_k z_{\omega,r,k},
    && r \in \mathcal{R},
    \label{eq:battery-size} \\
    & z_{\omega,r,k} \in \{0,1\},
    && r \in \mathcal{R},\ k \in \mathcal{K}.
    \label{eq:z-domain}
\end{align}
Constraints~\eqref{eq:one-battery}--\eqref{eq:z-domain} select one common battery option for all buses assigned to each route and define the induced route battery capacity $B_{\omega,r}$.

\paragraph{Shared plug-in charging.}
Plug-in charger use is modeled over scenario time slots $\mathcal{T}_\omega$.
A visit $v$ can use plug-in charging only during $\mathcal{T}_v\subseteq\mathcal{T}_\omega$, and $s(v)$ maps the visit to its station.
The binary variable $x_{v\tau}$ records whether visit $v$ occupies a plug-in charger in time slot $\tau$; each occupied slot delivers up to $\kappa$ units of grid-side energy, and station $s$ has capacity $U_s$.
The plug-in charging constraints are
\begin{align}
    & 0
    \leq
    q_v^{\mathrm{P}}
    \leq
    \kappa
    \sum_{\tau \in \mathcal{T}_v} x_{v\tau},
    && r \in \mathcal{R},\ v \in \mathcal{V}^{\mathrm{P}}_{\omega,r},
    \label{eq:plugin-charge-limit} \\
    & \sum_{r \in \mathcal{R}}
    \sum_{\substack{v \in \mathcal{V}^{\mathrm{P}}_{\omega,r}:\\ s(v)=s,\ \tau\in\mathcal{T}_v}}
    x_{v\tau}
    \leq
    U_s,
    && s \in \mathcal{C},\ \tau \in \mathcal{T}_\omega,
    \label{eq:plugin-capacity} \\
    & x_{v\tau} \in \{0,1\},
    && r \in \mathcal{R},\ v \in \mathcal{V}^{\mathrm{P}}_{\omega,r},\ \tau \in \mathcal{T}_v.
    \label{eq:x-domain}
\end{align}
Constraints~\eqref{eq:plugin-charge-limit}--\eqref{eq:x-domain} define plug-in charging amounts, plug-in charger occupancy, and station-time capacity.
The shared plug-in charger-capacity constraint~\eqref{eq:plugin-capacity} provides the main cross-route coupling within a scenario.

\paragraph{Route-local operational feasibility.}
Route energy feasibility is written over station visits $\mathcal{V}_{\omega,r}$ and directed traversals $\mathcal{A}_{\omega,r}$.
The mappings $\ell(a)$, $\mathrm{tail}(a)$, and $\mathrm{head}(a)$ connect each traversal to its physical link and adjacent visits.
Traversal and depot energy requirements are represented by $d_a$, $d_m^0$, and $d_m^N$, while $\bar q$ and $\eta$ collect charging-opportunity limits and charging efficiencies.
The variables $e_v^{\mathrm{arr}}$ and $e_v^{\mathrm{dep}}$ are the battery energy states at arrival and departure of visit $v$.
The state-of-charge fractions are bounded by $\underline{\rho}$ and $\overline{\rho}$, and $v_m^0$ and $v_m^N$ denote the first and last visits of trip chain $m$; detailed notation is listed in Appendix~\ref{app:model-details}.
Trip-level energy requirements are exogenous and do not vary with the selected battery capacity.
Overnight depot charging is assumed to have sufficient power and capacity to restore all buses to the prescribed initial state of charge before the next representative day.
Event-level variables suppress scenario and route indices because each visit and traversal is already a scenario-route-specific copy.
Charging quantities are nonnegative continuous grid-side energy variables, and battery-energy states are continuous.
For each route $r\in\mathcal R$, the route-local operational constraints are
\begin{align}
    & 0
    \leq
    q_v^{\mathrm{S}}
    \leq
    \bar{q}^{\mathrm{S}}_v w_{s(v)}^{\mathrm{S}},
    && v \in \mathcal{V}^{\mathrm{S}}_{\omega,r},
    \label{eq:swc-charge-limit} \\
    & 0
    \leq
    q_a^{\mathrm{D}}
    \leq
    \bar{q}^{\mathrm{D}}_a w_{\ell(a)}^{\mathrm{D}},
    && a \in \mathcal{A}_{\omega,r}.
    \label{eq:dwc-charge-limit} \\
    & \underline{\rho} B_{\omega,r}
    \leq
    e^{\mathrm{arr}}_v
    \leq
    \overline{\rho} B_{\omega,r},
    && v \in \mathcal{V}_{\omega,r},
    \label{eq:soc-arr-bound} \\
    & \underline{\rho} B_{\omega,r}
    \leq
    e^{\mathrm{dep}}_v
    \leq
    \overline{\rho} B_{\omega,r},
    && v \in \mathcal{V}_{\omega,r},
    \label{eq:soc-dep-bound} \\
    & e^{\mathrm{dep}}_v
    =
    e^{\mathrm{arr}}_v
    +
    \eta^{\mathrm{P}}
    q_v^{\mathrm{P}},
    && v \in \mathcal{V}^{\mathrm{P}}_{\omega,r},
    \label{eq:station-energy-balance} \\
    & e^{\mathrm{dep}}_v
    =
    e^{\mathrm{arr}}_v
    +
    \eta^{\mathrm{S}}
    q_v^{\mathrm{S}},
    && v \in \mathcal{V}^{\mathrm{S}}_{\omega,r}, \notag \\
    & e^{\mathrm{arr}}_{\mathrm{head}(a)}
    =
    e^{\mathrm{dep}}_{\mathrm{tail}(a)}
    -
    d_a
    +
    \eta^{\mathrm{D}}
    q_a^{\mathrm{D}},
    && a \in \mathcal{A}_{\omega,r},
    \label{eq:link-energy-balance} \\
    & e^{\mathrm{arr}}_{v_m^{0}}
    =
    \overline{\rho} B_{\omega,r}
    -
    d_m^{0},
    && m \in \mathcal{M}_{\omega,r},
    \label{eq:initial-boundary-energy} \\
    & e^{\mathrm{dep}}_{v_m^{N}}
    -
    d_m^{N}
    \geq
    \underline{\rho} B_{\omega,r},
    && m \in \mathcal{M}_{\omega,r},
    \label{eq:return-boundary-energy} \\
    & \eta^{\mathrm{P}} q_m^{\mathrm{N}}
    =
    \overline{\rho} B_{\omega,r}
    -
    e^{\mathrm{dep}}_{v_m^{N}}
    +
    d_m^{N},
    && m \in \mathcal{M}_{\omega,r}.
    \label{eq:overnight-charge-balance}
\end{align}
Constraints~\eqref{eq:swc-charge-limit}--\eqref{eq:dwc-charge-limit} link route operations to first-stage wireless-facility decisions, and constraints~\eqref{eq:soc-arr-bound}--\eqref{eq:soc-dep-bound} impose arrival and departure state-of-charge bounds.
Constraints~\eqref{eq:station-energy-balance}--\eqref{eq:link-energy-balance} update the energy state through station charging and directed traversals, while constraints~\eqref{eq:initial-boundary-energy}--\eqref{eq:overnight-charge-balance} close each trip chain through depot movements and overnight recharge.

\subsection{Recourse structure}
\label{subsec:recourse-structure}

For a fixed facility layout $w$, the expected recourse term is separable across scenarios. This yields a standard integer L-shaped reference in which each scenario is evaluated as a monolithic mixed-integer recourse problem over $(z_\omega,x_\omega,y_\omega)$. We refer to this direct decomposition as Natural ILS.
The recourse problem also exhibits conditional route locality. Constraints~\eqref{eq:swc-charge-limit}--\eqref{eq:overnight-charge-balance} are imposed independently by route, whereas the shared plug-in charger-capacity constraint~\eqref{eq:plugin-capacity} couples plug-in occupancy across routes at each station--time pair.
Finally, the formulation induces a facility--battery resource order. Under the model specification that trip-level energy consumption is independent of the selected battery capacity, additional wireless facilities and larger route batteries preserve operational feasibility under resource-dominating configurations and cannot increase the minimum charging cost.

These properties define a layered recourse structure comprising scenario-dependent discrete battery configurations, conditionally route-local operations, and limited shared-capacity coupling. Section~\ref{sec:method} develops an exact decomposition method tailored to this structure.
